% Generated draft: requires mathematical and source review.
\documentclass[11pt]{article}
\usepackage[margin=25mm]{geometry}
\usepackage{fontspec}
\setmainfont{Latin Modern Roman}
\usepackage{amsmath,amssymb,mathtools,unicode-math}
\setmathfont{Latin Modern Math}
\usepackage{enumitem,needspace,xurl,hyperref}
\setlength{\parindent}{0pt}
\setlength{\parskip}{5pt}
\setlength{\emergencystretch}{3em}
\newcommand{\N}{\mathbb{N}}
\newcommand{\Z}{\mathbb{Z}}
\newcommand{\Q}{\mathbb{Q}}
\newcommand{\R}{\mathbb{R}}
\newcommand{\C}{\mathbb{C}}
\newcommand{\F}{\mathbb{F}}
\newcommand{\A}{\mathbb{A}}
\newcommand{\PP}{\mathbb{P}}
\newcommand{\E}{\mathbb{E}}
\newcommand{\eps}{\varepsilon}
\newcommand{\dd}{\,\mathrm{d}}
\begin{document}
\section*{UnsolvedMath ID: 6900002}
\textbf{Source problem code:} catalogue entry
\subsection{Definitions and mathematical statement}
The supplied question asks for all different types of strata for ten points, in the planar tensegrity-configuration setting described by the research summary. The source excerpt does not provide the definition of the configuration space or its stratification.

For context, for a graph \(G\) on labeled vertices \(1,\ldots,10\) and positions \(p=(p_1,\ldots,p_{10})\in(\mathbb R^2)^{10}\), an equilibrium stress is an assignment \(\omega_{ij}=\omega_{ji}\in\mathbb R\) to its edges satisfying
\[
 \sum_{j:\{i,j\}\in E(G)}\omega_{ij}(p_j-p_i)=0
 \qquad(1\le i\le10).
\]
These assignments form a linear space \(W(G,p)\). Tensegrity stratifications can track more than its dimension, for instance sign data or sign-preserving equivalence of fibers. The question is to classify every stratum type in the source's ten-point space, not merely possible stress dimensions or codimension-one strata.

The graph, whether all pairs are edges, labeled-versus-unlabeled convention, admissible point coincidences/degeneracies, any geometric quotient, and precise equivalence defining a stratum ``type'' are missing. The displayed stress space gives standard context without choosing those absent ingredients. In particular a rank stratification alone is not asserted to be the intended tensegrity stratification.

The supplied summary quotes the source's progress: codimension-one strata for \(n\le8\) were classified by Doray and collaborators (2010), and Karpenkov (2017) gave an approach to every stratum for \(n=9\); it identifies \(n=10\) as the next case containing unknown strata. ``An approach'' is not silently upgraded to a complete lower-dimensional classification. The record reports the ten-point problem open, without an independently verified modern resolution.
\subsection{Short English statement}
Classify every stratum type in the tensegrity configuration space for ten planar points. The source's graph, admissible configurations, and definition of type are missing; the cited smaller-case progress concerns codimension-one classification and an approach to the nine-point case.
\subsection{Sources}
\begin{itemize}
\item Local registry: \url{data/unsolvedmath/problems.json}, permanent UnsolvedMath ID \(6900002\), source code catalogue entry; supplied statement, background, and research summary, updated 2026-08-06.
\item Oleg Karpenkov, \emph{Open Problems on Configuration Spaces of Tensegrities} (2020), Problem 2, PDF page 3; progress discussed in section 3: \url{https://livrepository.liverpool.ac.uk/3022838/1/stratum.pdf}. Supplied normalized PDF-text extraction accessed 2026-07-29.
\item Doray and collaborators (2010), cited by the supplied summary for codimension-one strata with \(n\le8\); no full title or bibliographic details are supplied.
\item Karpenkov (2017), arXiv:1512.02563, cited for the nine-point approach; no full title is supplied in the record.
\end{itemize}
\end{document}
